\documentclass[11pt]{article}

\usepackage[margin=1in]{geometry}
\usepackage{amsmath, amssymb}
\usepackage{enumitem}
\usepackage{hyperref}
\hypersetup{hidelinks}

\title{Scaling Book Exercises: Section 1}
\author{Lucas Yan}
\date{}

\newcommand{\AI}{\mathrm{AI}}
\newcommand{\tmath}{T_{\mathrm{math}}}
\newcommand{\tcomm}{T_{\mathrm{comm}}}

\begin{document}
\maketitle

\noindent
\textbf{Source scan:} \texttt{../handwritten/section\_01\_handwritten.pdf}

\section*{Exercise 1}

Consider the matrix multiplication
\[
X \in \mathbb{R}^{B \times D}, \qquad
Y \in \mathbb{R}^{D \times F}, \qquad
Z = XY \in \mathbb{R}^{B \times F}.
\]

\begin{enumerate}[label=(\alph*)]
  \item Loading the inputs costs
  \[
  BD + DF
  \]
  bytes, and writing the output costs
  \[
  BF
  \]
  bytes.

  \item The total number of floating-point operations is
  \[
  2BDF.
  \]

  \item The arithmetic intensity is therefore
  \[
  \AI = \frac{2BDF}{BD + DF + BF}.
  \]

  \item Using a tensor-core throughput of $3.94 \times 10^{14}$ FLOPs/s and communication bandwidth of $8.2 \times 10^{11}$ bytes/s,
  \[
  \tmath = \frac{2BDF}{3.94 \times 10^{14}},
  \qquad
  \tcomm = \frac{BD + DF + BF}{8.2 \times 10^{11}}.
  \]
  A reasonable upper bound on runtime is
  \[
  \tmath + \tcomm,
  \]
  and a reasonable lower bound is
  \[
  \max(\tmath, \tcomm).
  \]
\end{enumerate}

\section*{Exercise 2}

We solve for the batch size at which compute time and communication time are equal. Assume
\[
B \ll D, \qquad B \ll F.
\]
Using the arithmetic intensity threshold from the notes,
\[
\frac{2BDF}{DF} = 2B > 240.
\]
Therefore,
\[
B > 120.
\]

\section*{Exercise 3}

For a fixed batch size $B = 100$, the total bytes moved are approximately
\[
2BD + 2BF + DF.
\]
When $D = F$, this becomes
\[
4BD + D^2.
\]
The arithmetic intensity is
\[
\AI
  \approx
  \frac{2BD^2}{2BD + 2BF + DF}
  =
  \frac{2BD^2}{4BD + D^2}
  =
  \frac{2BD}{4B + D}.
\]
With $B = 100$,
\[
\AI = \frac{200D}{400 + D}.
\]

For $D = F = 4096$,
\[
\AI = \frac{200 \cdot 4096}{400 + 4096} \approx 182.
\]

For $D = F = 1024$,
\[
\AI = \frac{200 \cdot 1024}{400 + 1024} \approx 143.
\]

\section*{Exercise 4}

Arithmetic intensity is
\[
\AI = \frac{\mathrm{FLOPs}}{\mathrm{communication}}.
\]
For the matrix multiplication
\[
[B,D] @ [D,F],
\]
the FLOP count is
\[
2BDF.
\]
The byte count is
\[
BD + BDF + BF.
\]
Thus
\[
\AI = \frac{2BDF}{BD + BDF + BF}.
\]

Assuming $B \ll D$ and $B \ll F$, the dominant term in the denominator is $BDF$, so
\[
\AI \approx \frac{2BDF}{BDF} = 2.
\]
This is constant with respect to the matrix sizes under these assumptions.

\section*{Exercise 5}

The critical batch size is computed from the tensor-core throughput and bandwidth:
\[
\text{tensor core throughput} = 1979 \text{ teraFLOPs},
\qquad
\text{bandwidth} = 3.35 \text{ terabytes/s}.
\]
Accounting for multiply-add as two FLOPs,
\[
B_{\mathrm{critical}}
  =
  \frac{1979/2}{3.35}
  \approx 295.37.
\]
Thus,
\[
B_{\mathrm{critical}} \approx 295.
\]

\section*{Verification Notes}

The formulas above are converted from the handwritten scan. Before final submission, verify the following against the exercise text:

\begin{itemize}
  \item Exercise 4's byte-count expression, since the handwritten denominator appears as $BD + BDF + BF$.
  \item Whether the bandwidth constants should be expressed in bytes/s or words/s in the final writeup.
\end{itemize}

\end{document}
